\section{Introduction}
\label{sec:introduction}

For a finite graph $G$, a plane drawing places its vertices at distinct
points and represents each edge by a simple continuous arc whose interior
avoids the vertices. A meeting of two edge interiors is a proper crossing
if some open neighborhood of the point has a homeomorphism onto an open
disk taking the point to its center and the two local edge traces onto its
horizontal and vertical diameters. We require that edge interiors meet at
finitely many proper crossings, with no triple crossings, tangencies, or
common subarcs.
The number $c(D)$ of a drawing $D$ counts crossing points: if two edges
cross several times, each crossing is counted.
The ordinary crossing number is
\[
  \crn(G)=\min_D c(D).
\]
There are no restrictions on vertex positions or on edge routes.

Let $K_{m,n}$ be the complete bipartite graph with vertex classes
$\{a_1,\ldots,a_m\}$ and $\{b_1,\ldots,b_n\}$. Put
\begin{equation}
\label{eq:dr}
 d_r=\left\lfloor\frac r2\right\rfloor
     \left\lfloor\frac{r-1}{2}\right\rfloor
 =\binom{\lfloor r/2\rfloor}{2}
  +\binom{\lceil r/2\rceil}{2}.
\end{equation}
\begin{theorem}
\label{thm:main}
For all positive integers $m,n$,
\[
 \crn(K_{m,n})=d_m d_n.
\]
\end{theorem}

The distinction between crossing points, crossing pairs, and parity-based
crossing counts is important; we use the ordinary point-count convention
throughout~\cite{PachToth2000,FulekEtAl2012}.

Tur\'an described the problem's origin in rail connections between kilns
and storage yards at a brick factory near Budapest in
1944~\cite[pp.~8--9]{Turan1977}.
Zarankiewicz supplied the balanced-axis construction, but the published
argument for its optimality contained a gap
\cite{Zarankiewicz1955,Turan1977,ErdosGuy1973}.
Kleitman established the exact formula when one vertex class has at
most six vertices~\cite{Kleitman1970}.
His permutation argument already used the local cyclic orders of incident
edges to bound crossings between two stars
\cite[Section~8, equation~(17)]{Kleitman1970}.
Woodall developed this approach through cyclic-order graphs and verified the cases
$K_{7,7}$ and $K_{7,9}$ by computation~\cite{Woodall1993}.
The rotation constraints later supported semidefinite lower bounds
\cite{deKlerkEtAl2006,deKlerkPasechnikSchrijver2007}, refined for fixed
part sizes by Brosch and Polak~\cite{BroschPolak2024}.
For balanced graphs, Balogh, Lidick\'y, Norin, Pfender, Salazar, and
Spiro obtained
\[
 \operatorname{cr}(K_{n,n})\ge0.9118\,d_n^2+o(n^4)
\]
using Razborov's flag-algebra method~\cite{Razborov2007,BaloghEtAl2023}.
Hebbar and Mangam have also published claims of both the complete-graph
and complete-bipartite formulas~\cite{HebbarMangam2018}.
Theorem~\ref{thm:main} resolves the Zarankiewicz crossing-number
conjecture positively through the self-contained argument below.

Tutte developed an algebraic approach to crossing
numbers~\cite{Tutte1970}. Our proof combines signed intersections of edge
interiors with a contribution determined by the incident-edge order at a
common endpoint. The resulting cycle pairing separates the topological
information from the algebraic estimates.

The proof separates the geometry of a drawing from two linear-algebraic
statements. In Section~\ref{sec:subspaces}, we prove that subspaces
$U_i,V_i$ of an $N$-dimensional
real vector space satisfying $U_i\cap V_i=0$ obey
\begin{equation}
\label{eq:intro-subspaces}
 \sum_{i,j=1}^m\dim(U_i\cap V_j)
 \le \left\lfloor\frac{m^2}{4}\right\rfloor N.
\end{equation}
The argument uses interpolation with subspace constraints at finitely many
nodes and counts the zeros of polynomial minors. Related determinant
multiplicity arguments control sums of intersection dimensions in the
theory of subspace designs~\cite{GuruswamiKopparty2016}.
Here the degree filtration is adapted to arbitrary subspaces $U_i$;
diagonal disjointness bounds the rank at the interpolation nodes.
Applied to graphs of linear maps, the resulting inequality gives a lower
bound on the sum of the ranks of differences $X_i-Y_j$ of square
matrices, provided every diagonal difference $X_i-Y_i$ is invertible.
Both formulations are independent of graph drawings.

In Section~\ref{sec:detector}, for odd $n=2s+1\ge3$, we construct a linear map from $n$-by-$n$
matrices to $s^2$-by-$s^2$ matrices. It annihilates row-plus-column terms
and diagonal entries and sends each individual matrix entry to a matrix of
rank at most one. For a linear order $\prec$ on $\{1,\ldots,n\}$, let
$S_\prec$ be half its comparison-sign matrix: its $(p,q)$ entry is
$\tfrac12$ if $p\prec q$, $-\tfrac12$ if $q\prec p$, and $0$ if $p=q$.
The map sends every $S_\prec$ to an invertible matrix.
The construction uses the classical barycentric weights for Lagrange interpolation
\cite[Sections~2--3]{BerrutTrefethen2004} and bivariate monomial evaluations;
its invertibility reduces to interpolation on two complementary triangles
of an evaluation grid.

To connect these statements, normalize a minimum drawing $D$ so its
edges are polygonal and adjacent edges do not cross. Orient edges from
the first vertex class to the second. At each vertex, read the counterclockwise
spoke order from a chosen gap. Form a matrix $J$ indexed by these
edges. Its entry for two distinct edges is their signed crossing count,
plus half their comparison sign in the incident order if they share an
endpoint; diagonal entries are zero. The block $J^{ik}$
compares the edges incident with $a_i$ to those incident with $a_k$.
The vanishing of algebraic intersections of closed plane curves makes
these blocks additive modulo row-plus-column terms
(Section~\ref{sec:intersections}).
Applying the linear map produces
the family $X_i-Y_j$: diagonal blocks have full rank, and off-diagonal
ranks are bounded by the corresponding crossing counts.
Corollary~\ref{cor:rank-sum} bounds the ordered sum of off-diagonal
ranks below by $2d_m d_n$. The corresponding sum of crossing counts
counts each crossing twice, so the rank sum is at most $2c(D)$.
This gives $c(D)\ge d_m d_n$.
Section~\ref{sec:formula} completes the proof: deleting vertices gives
the even case, and the classical balanced-axis drawing gives the
matching upper bound.

Every ingredient of the lower-bound argument is proved below.
For completeness, Appendix~\ref{sec:normalization} supplies the drawing
reductions for the continuous-arc definition used here.
